\section{Detector and Evaluation Setup}\label{sec:setup}

\emph{CIC-IoT2023.} We use ten captures~\cite{neto2023ciciot2023}: Benign, DDoS-SYN, DoS-HTTP, Recon-PortScan, MITM-ArpSpoofing, DictionaryBruteForce, SqlInjection, XSS, CommandInjection and Mirai-udpplain (about 5.17\,GB). Files over 300\,MB are trimmed to their first $10^6$ packets. The captures map to eight classes: Benign, BruteForce, DDoS-SYN, DoS-HTTP, MITM-ARP, Mirai, Recon and Web (SQLi, XSS and command injection).

\emph{IoT-23.} We use ten IoT-23 scenarios~\cite{garcia2020iot23}: seven malware captures (Hide\&Seek, Muhstik, Hakai, two Torii, Mirai, Trojan) and three benign captures of real, uninfected IoT devices (Philips Hue bridge, Amazon Echo, Somfy smart door lock; of 7-1 only its Somfy-01 part). The Okiru captures (all at least 1\,GB) are left out. Labels are the dataset's own Zeek \texttt{conn.log.labeled}: each flow is joined to the Zeek connection with the same unordered 5-tuple and the nearest start within $\pm2$\,s, otherwise to the Zeek connection whose span contains the flow's start. Of 756,306 flows, 99.99\% match; the 38 unmatched flows are dropped. The six classes are Benign, PortScan, C\&C, Attack (only in scenario 3-1), DDoS and FileDownload. Of the flows, 35.2\% overlap more than one Zeek row, but only 270 (0.04\%) overlap rows of more than one class, so the labels agree at flow level. Class supports differ, because flows and Zeek connections are different units: 14,394 Zeek DDoS rows on 211 5-tuples become 211 NFStream flows, one of which covers 14,159 rows.

\emph{Flows and features.} Flows are built with NFStream~6.6.0~\cite{aouini2022nfstream} (\texttt{statistical\_analysis=True}, idle timeout 15\,s, active timeout 120\,s), with identical settings for both datasets and on the live gateway. The model uses 54 statistical features (packet counts and sizes, inter-arrival times, durations, TCP flag counts per direction). Identifiers (addresses, ports, protocol, timestamps) are excluded, because they let a model memorise the capture~\cite{arp2022dos}. CIC-IoT2023 yields 857,206 labelled flows, 511,527 of them DDoS-SYN (a consequence of trimming).

\emph{Detector.} The detector is a scikit-learn~\cite{pedregosa2011sklearn} decision tree with balanced class weights. On the CIC-IoT2023 temporal split it scores 0.731, close to a random forest (0.737, at about 94$\times$ the size and 58$\times$ the single-flow latency) and histogram gradient boosting (0.714, Benign FPR 7.6\%). Limiting depth and leaf size on validation raised the validation Benign FPR (from 1.3\% to 5.7\% at depth 16, leaf 10), so the deployed tree (depth 32, leaf 1; 1,767 nodes, 222\,KiB) is effectively unconstrained; it scores 0.725 (Benign FPR 3.8\%). NSGA-II feature selection~\cite{deb2002nsga2,blank2020pymoo} kept 14 of 54 features at no measurable cost (0.723) but barely shrank the tree (211\,KiB). Section~\ref{sec:splits} reports the baseline tree (0.731, Benign FPR 3.9\%); all other results use the deployed tree, retrained with the same settings on IoT-23.

\emph{Splits.} The \emph{temporal} split puts the first 70\% of each capture's flows, by start time, into training (on CIC-IoT2023, 600,050 training and 257,156 test flows). The \emph{stratified random} split draws 70\% at random, stratified by class. Every design choice (tree limits, feature subset, blocking policy) is made on a validation set, the last 20\% of each capture's training period.

\emph{Metrics.} We report per-class F1 with test support, macro-F1 over all classes, and the Benign false-positive rate (FPR), the share of Benign test flows predicted as any attack. Flows from one host pair are not independent, so 95\% confidence intervals come from a \emph{host-pair bootstrap}: 1,000 resamples of whole (capture, host pair) clusters. A flow-level bootstrap~\cite{efron1994bootstrap} gives far too narrow intervals: [0.712, 0.742] for the temporal decision tree, against [0.577, 0.823] with host pairs. Two splits are compared by a \emph{paired} bootstrap over the union of both test sets' (capture, unordered host pair) clusters (CIC-IoT2023: 4,641 from 3,465 temporal and 3,858 random; IoT-23: 348,350): each resample draws union clusters with replacement and applies the same draw to both splits (a cluster absent from one test set adds nothing to it), and F1 is recomputed from the summed confusion matrices. All reported intervals use host-pair clusters.
